\documentclass{article}
\usepackage{graphicx} % Required for inserting images
\usepackage{fancyhdr}
\usepackage{amsmath}
\usepackage{float}



\begin{document}
\pagestyle{fancy}
\fancyhead[RO]{Mehmet Çağatay Kurmaç}
\fancyhead[LO]{Target Tracking Exercise 3}

\section*{a)}

    The true trajectory of the target is shown in Figure \ref{fig:1a}. 
    \begin{figure}[H]
        \centering
        \includegraphics[width=0.8\textwidth]{figures/true_traj.png}
        \caption{The plot of the true trajectory.}
        \label{fig:1a}
    \end{figure}

\section*{b)}

    The noisy measurements of the target are shown in Figure \ref{fig:1b}.

    \begin{figure}[H]
        \centering
        \includegraphics[width=0.8\textwidth]{figures/noisy_meas.png}
        \caption{The plot of the target and the measurements.}
        \label{fig:1b}
    \end{figure}

\section*{c)}

    For Q1, the estimated trajectory with gates are given in Figure \ref{fig:1c}. The estimation errors with respect to time are given
    in Figure \ref{fig:1d}. The gate sizes are given in Figure \ref{fig:1e}.

    \begin{figure}[H]
        \centering
        \includegraphics[width=0.8\textwidth]{figures/q1_traj.png}
        \caption{The plot of the target, measurements, state estimates and gates for Q1.}
        \label{fig:1c}
    \end{figure}

    \begin{figure}[H]
        \centering
        \includegraphics[width=0.8\textwidth]{figures/q1_est_err.png}
        \caption{The plot of the estimation errors with respect to time for Q1.}
        \label{fig:1d}
    \end{figure}

    \begin{figure}[H]
        \centering
        \includegraphics[width=0.8\textwidth]{figures/q1_gate_size.png}
        \caption{The plot of the gate sizes with respect to time for Q1.}
        \label{fig:1e}
    \end{figure}

    For Q2, the estimated trajectory with gates are given in Figure \ref{fig:1f}. The estimation errors with respect to time are given
    in Figure \ref{fig:1g}. The gate sizes are given in Figure \ref{fig:1h}.

    \begin{figure}[H]
        \centering
        \includegraphics[width=0.8\textwidth]{figures/q2_traj.png}
        \caption{The plot of the target, measurements, state estimates and gates for Q2.}
        \label{fig:1f}
    \end{figure}

    \begin{figure}[H]
        \centering
        \includegraphics[width=0.8\textwidth]{figures/q2_est_err.png}
        \caption{The plot of the estimation errors with respect to time for Q2.}
        \label{fig:1g}
    \end{figure}

    \begin{figure}[H]
        \centering
        \includegraphics[width=0.8\textwidth]{figures/q2_gate_size.png}
        \caption{The plot of the gate sizes with respect to time for Q2.}
        \label{fig:1h}
    \end{figure}

        For Q3, the estimated trajectory with gates are given in Figure \ref{fig:1i}. The estimation errors with respect to time are given
    in Figure \ref{fig:1j}. The gate sizes are given in Figure \ref{fig:1k}.

    \begin{figure}[H]
        \centering
        \includegraphics[width=0.8\textwidth]{figures/q3_traj.png}
        \caption{The plot of the target, measurements, state estimates and gates for Q3.}
        \label{fig:1i}
    \end{figure}

    \begin{figure}[H]
        \centering
        \includegraphics[width=0.8\textwidth]{figures/q3_est_err.png}
        \caption{The plot of the estimation errors with respect to time for Q3.}
        \label{fig:1j}
    \end{figure}

    \begin{figure}[H]
        \centering
        \includegraphics[width=0.8\textwidth]{figures/q3_gate_size.png}
        \caption{The plot of the gate sizes with respect to time for Q3.}
        \label{fig:1k}
    \end{figure}

\section*{d)}

    The IMM Filter for given transition probability matrices are implemented with CV Kalman Filters given for Q1 and Q3. The resulting
    estimated trajectories with gates, true trajectory and measurements are given in Figure \ref{fig:1l}. The estimation errors with respect 
    to time are given in Figure \ref{fig:1m}. The gate sizes are given in Figure \ref{fig:1n}. The model probabilities with respect to time
    are given in Figure \ref{fig:1o}.

    \begin{figure}[H]
        \centering
        \includegraphics[width=0.8\textwidth]{figures/d_imm_traj.png}
        \caption{The plot of the target, measurements, state estimates and gates for IMM Filter.}
        \label{fig:1l}
    \end{figure}

    \begin{figure}[H]
        \centering
        \includegraphics[width=0.8\textwidth]{figures/d_imm_est_err.png}
        \caption{The plot of the estimation errors with respect to time for IMM Filter.}
        \label{fig:1m}
    \end{figure}

    \begin{figure}[H]
        \centering
        \includegraphics[width=0.8\textwidth]{figures/d_imm_gate_size.png}
        \caption{The plot of the gate sizes with respect to time for IMM Filter.}
        \label{fig:1n}
    \end{figure}

    \begin{figure}[H]
        \centering
        \includegraphics[width=0.8\textwidth]{figures/d_imm_model_prob.png}
        \caption{The plot of the model probabilities with respect to time for IMM Filter.}
        \label{fig:1o}
    \end{figure}

\section*{e)}

    Two IMM filters with the altered transition probability matrices are implemented.

    For TPM 1 (0.999,0.001), the resulting estimated trajectories with gates, true trajectory and measurements are given in Figure \ref{fig:1p}. The estimation errors with respect 
    to time are given in Figure \ref{fig:1q}. The gate sizes are given in Figure \ref{fig:1r}. The model probabilities with respect to time
    are given in Figure \ref{fig:1s}.

    \begin{figure}[H]
        \centering
        \includegraphics[width=0.8\textwidth]{figures/e_imm_1_traj.png}
        \caption{The plot of the target, measurements, state estimates and gates for IMM Filter with TPM 1.}
        \label{fig:1p}
    \end{figure}

    \begin{figure}[H]
        \centering
        \includegraphics[width=0.8\textwidth]{figures/e_imm_1_est_err.png}
        \caption{The plot of the estimation errors with respect to time for IMM Filter with TPM 1.}
        \label{fig:1q}
    \end{figure}

    \begin{figure}[H]
        \centering
        \includegraphics[width=0.8\textwidth]{figures/e_imm_1_gate_size.png}
        \caption{The plot of the gate sizes with respect to time for IMM Filter with TPM 1.}
        \label{fig:1r}
    \end{figure}

    \begin{figure}[H]
        \centering
        \includegraphics[width=0.8\textwidth]{figures/e_imm_1_model_prob.png}
        \caption{The plot of the model probabilities with respect to time for IMM Filter with TPM 1.}
        \label{fig:1s}
    \end{figure}
        
    For TPM 2 (0.5,0.5), the resulting estimated trajectories with gates, true trajectory and measurements are given in Figure \ref{fig:1t}. The estimation errors with respect 
    to time are given in Figure \ref{fig:1u}. The gate sizes are given in Figure \ref{fig:1v}. The model probabilities with respect to time
    are given in Figure \ref{fig:1w}.

    \begin{figure}[H]
        \centering
        \includegraphics[width=0.8\textwidth]{figures/e_imm_2_traj.png}
        \caption{The plot of the target, measurements, state estimates and gates for IMM Filter with TPM 2.}
        \label{fig:1t}
    \end{figure}

    \begin{figure}[H]
        \centering
        \includegraphics[width=0.8\textwidth]{figures/e_imm_2_est_err.png}
        \caption{The plot of the estimation errors with respect to time for IMM Filter with TPM 2.}
        \label{fig:1u}
    \end{figure}

    \begin{figure}[H]
        \centering
        \includegraphics[width=0.8\textwidth]{figures/e_imm_2_gate_size.png}
        \caption{The plot of the gate sizes with respect to time for IMM Filter with TPM 2.}
        \label{fig:1v}
    \end{figure}

    \begin{figure}[H]
        \centering
        \includegraphics[width=0.8\textwidth]{figures/e_imm_2_model_prob.png}
        \caption{The plot of the model probabilities with respect to time for IMM Filter with TPM 2.}
        \label{fig:1w}
    \end{figure}

    As we can see from the results, If the transition probabilities are set to be more certain (TPM 1), the IMM filter performs better
    in terms of estimation error and gate sizes adapt to the altering dynamics as expected. On the other hand, if the transition probabilities
    are set to be more uncertain (TPM 2), the IMM filter performs worse in terms of estimation error and gate sizes do not adapt at all. The
    model probabilities change trivially and adaptation cannot be observed.

    This is expected because when we are calculating the estimates, the models act as hypotheses. If the transitions are allowed to be made 
    loosely, the hypothesis aspect becomes trivial and representation ability perishes.
\end{document}